\section{Metodologia}\label{metodologia}

\subsection{Desenho da pesquisa}\label{desenho-da-pesquisa}

O estudo é empírico, quantitativo e preditivo, com uma camada
interpretativa; seu objetivo é medir informação preditiva, não estimar
efeitos causais \cite{shmueli2010explain}. A fase de
\textbf{desenvolvimento} (exploratória) inclui as auditorias do dataset,
a definição dos alvos e das visões de atributos, a validação agrupada e
a seleção de modelo. A fase de \textbf{confirmação} avalia, uma única
vez, os modelos congelados em um holdout novo, com critérios de sucesso
fixados antes. O protocolo completo, inclusive as decisões tomadas
depois de ver resultados, está em \texttt{paper/protocol.md}.

\subsection{Dados e unidade de
análise}\label{dados-e-unidade-de-anuxe1lise}

A fonte é o corpus GitSkills \cite{gitskills2027}, revisão
\texttt{ebab17454a7c236f8f26b183567f1a126f42e3f8}: 3.797.117 arquivos
\texttt{SKILL.md} de 282.200 repositórios (julho de 2026), agrupados em
1.877.981 conteúdos distintos. A unidade é o conteúdo
(\texttt{content\_hash}). Todas as ocorrências (repositório e caminho)
são mantidas, mas somente um \textbf{representante} por conteúdo recebe
texto completo, metadados do repositório, pasta e histórico.

O desenvolvimento usa 200.000 conteúdos escolhidos por
\texttt{md5(content\_hash\ \textbar{}\textbar{}\ \textquotesingle{}:20261009\textquotesingle{})},
divididos por repositório em 160.000 de desenvolvimento e 40.000 de
teste. Condicionamos a \texttt{repository\_count\ \textgreater{}=\ 2}:
25.918 conteúdos de 7.128 donos no desenvolvimento. Na população
completa há 306.088 conteúdos nessa condição, dos quais 230.355 sob dois
ou mais donos. A partição de teste original foi consultada na primeira
versão do estudo e não sustenta afirmações confirmatórias. A data de
referência para idade e recência é 20/07/2026.

\subsection{Regra do representante e
auditoria}\label{regra-do-representante-e-auditoria}

A documentação do GitSkills indica que o representante prefere
\texttt{.claude/skills/}, com desempate determinístico
\cite{gitskills2027}. Recuperamos a regra a partir de todas as
ocorrências brutas: ordenar por classe de localização
(\texttt{canonical}, \texttt{skills-dir}, \texttt{other}) e, no empate,
pelo menor par \texttt{(repo\_full\_name,\ path)}. Ela reproduz o
representante em 97,2\% dos 306.088 conteúdos com ao menos dois
repositórios, e a classe do representante coincide com a melhor classe
disponível em 100\% deles. A auditoria compara, por faixa de número de
repositórios, a fração de conteúdos cujo representante é
\texttt{canonical} com a fração das cópias que de fato estão nessa
classe, para saber se \texttt{location\_class} e
\texttt{history\_fetched} (histórico só coletado para locais padrão e
uma amostra) podem ser preditores legítimos.

\subsection{Alvos}\label{alvos}

\begin{longtable}[]{@{}
  >{\raggedright\arraybackslash}p{(\linewidth - 4\tabcolsep) * \real{0.3333}}
  >{\raggedright\arraybackslash}p{(\linewidth - 4\tabcolsep) * \real{0.3333}}
  >{\raggedright\arraybackslash}p{(\linewidth - 4\tabcolsep) * \real{0.3333}}@{}}
\caption{Alvos. O dono é o prefixo de
\texttt{repo\_full\_name}.}\tabularnewline
\toprule\noalign{}
\begin{minipage}[b]{\linewidth}\raggedright
Papel
\end{minipage} & \begin{minipage}[b]{\linewidth}\raggedright
Alvo
\end{minipage} & \begin{minipage}[b]{\linewidth}\raggedright
Definição
\end{minipage} \\
\midrule\noalign{}
\endfirsthead
\toprule\noalign{}
\begin{minipage}[b]{\linewidth}\raggedright
Papel
\end{minipage} & \begin{minipage}[b]{\linewidth}\raggedright
Alvo
\end{minipage} & \begin{minipage}[b]{\linewidth}\raggedright
Definição
\end{minipage} \\
\midrule\noalign{}
\endhead
\bottomrule\noalign{}
\endlastfoot
Principal (classificação) & \texttt{cross\_owner} &
\texttt{owner\_count\ \textgreater{}=\ 2} \\
Secundário (regressão) & \texttt{log1p(owner\_count)} & donos distintos
entre todas as ocorrências \\
Sensibilidade & \texttt{log1p(repository\_count)} & alvo da primeira
versão \\
\end{longtable}

Contar donos remove a autocópia entre repositórios de uma mesma conta.
As contagens são limites inferiores, pois a busca de código do GitHub
indexa apenas a branch padrão, arquivos de até 384 KB e repositórios
recentes \cite{gitskills2027}.

\subsection{Atributos}\label{atributos}

\begin{longtable}[]{@{}
  >{\raggedright\arraybackslash}p{(\linewidth - 4\tabcolsep) * \real{0.3000}}
  >{\raggedleft\arraybackslash}p{(\linewidth - 4\tabcolsep) * \real{0.4000}}
  >{\raggedright\arraybackslash}p{(\linewidth - 4\tabcolsep) * \real{0.3000}}@{}}
\caption{Blocos de atributos. Os modelos principais usam os quatro
primeiros (52 colunas, visão \emph{limpa}); os atalhos entram só como
achado.}\tabularnewline
\toprule\noalign{}
\begin{minipage}[b]{\linewidth}\raggedright
Bloco
\end{minipage} & \begin{minipage}[b]{\linewidth}\raggedleft
n
\end{minipage} & \begin{minipage}[b]{\linewidth}\raggedright
Conteúdo
\end{minipage} \\
\midrule\noalign{}
\endfirsthead
\toprule\noalign{}
\begin{minipage}[b]{\linewidth}\raggedright
Bloco
\end{minipage} & \begin{minipage}[b]{\linewidth}\raggedleft
n
\end{minipage} & \begin{minipage}[b]{\linewidth}\raggedright
Conteúdo
\end{minipage} \\
\midrule\noalign{}
\endhead
\bottomrule\noalign{}
\endlastfoot
\texttt{skill} & 34 & estrutura do \texttt{SKILL.md} (tamanho, títulos,
listas, código, links, front matter, scripts, pasta) e nome do
arquivo \\
\texttt{repo} & 8 & \texttt{log1p} de estrelas e forks, fork, idade,
recência do push, dias ativos, linguagem, licença \\
\texttt{catalog} & 2 & \texttt{log1p} do número de skills no repositório
e indicador de pelo menos 10 \\
\texttt{path} & 8 & profundidade, raiz, \texttt{templates},
\texttt{plugins}, \texttt{vendor}, \texttt{docs}, \texttt{packages/src},
diretório de outra ferramenta \\
\emph{atalhos} & 4 & \texttt{location\_class},
\texttt{history\_fetched}, \texttt{commit\_count}, \texttt{.claude} no
caminho \\
\end{longtable}

Os atalhos dependem da regra do representante e ficam fora dos modelos
principais. Os atributos descrevem o representante, não todas as cópias;
agregá-los sobre todas as cópias dependeria das mesmas ocorrências que
definem o alvo e foi evitado. Imputação (mediana com indicador de
ausência), escala e codificação (one-hot, \texttt{min\_frequency=10})
são ajustadas dentro de cada dobra.

\subsection{Validação e limite por
dono}\label{validauxe7uxe3o-e-limite-por-dono}

A validação é agrupada pelo \textbf{dono do repositório representativo}:
3 repetições de 3 dobras (\texttt{StratifiedGroupKFold}, sementes
2027--2029, estratos por faixa de \texttt{owner\_count}), com as mesmas
partições para todos os modelos. Como poucas contas concentram a
amostra, aplicamos um \textbf{limite de 50 conteúdos por dono} (ordem de
\texttt{content\_hash}), restando 17.733 conteúdos. O valor 50 foi
escolhido \emph{depois} de observar a sensibilidade do AUC ao tratamento
de donos concentrados (Seção 4.4); é, portanto, uma decisão dependente
dos dados. Por dobra, medimos também a fração de linhas de validação
cujo \texttt{name} do front matter aparece no treino e as métricas no
subconjunto sem sobreposição (um agrupamento conjunto por dono e
\texttt{name} forma um componente com 64\% das linhas e foi descartado).

\subsection{Modelos, baselines e
tuning}\label{modelos-baselines-e-tuning}

O modelo principal é o Random Forest (100 árvores em regressão, 300 em
classificação; \texttt{max\_depth=12}, \texttt{min\_samples\_leaf=3},
\texttt{max\_features=sqrt}). As referências lineares são Ridge
(\(\alpha=10\)) e regressão logística (\(C=1\)); as baselines são a
prevalência e a média e a mediana de \texttt{log1p}. O tuning é
aninhado: seis famílias (linear, KNN, árvore, Random Forest, Gradient
Boosting, HistGradientBoosting) em duas variantes (hiperparâmetros
padrão e buscados por \texttt{RandomizedSearchCV}) sobre a visão limpa.
A validação externa usa 2 repetições de 3 dobras agrupadas por dono; a
busca interna, 3 dobras agrupadas por dono e 6 candidatas. A combinação
com melhor média externa é congelada (se for ``buscada'', a busca é
repetida em todo o desenvolvimento). O protocolo previa 3 repetições e
12 candidatas; reduzimos por custo computacional.

\subsection{Métricas, inferência e
interpretação}\label{muxe9tricas-inferuxeancia-e-interpretauxe7uxe3o}

As métricas principais são o AUC-ROC (classificação) e o RMSLE
(regressão); as secundárias, AUC-PR, Brier, log-loss, MAE e Spearman. As
comparações são pareadas por (repetição, dobra), com IC95\% e teste t
com \textbf{correção de Nadeau--Bengio} para amostras repetidas; o
Wilcoxon pareado é complementar. Cinco hipóteses pré-especificadas
formam uma família corrigida por Holm (\(\alpha=0{,}05\)); as demais
comparações são exploratórias.

\begin{longtable}[]{@{}
  >{\raggedright\arraybackslash}p{(\linewidth - 4\tabcolsep) * \real{0.3333}}
  >{\raggedright\arraybackslash}p{(\linewidth - 4\tabcolsep) * \real{0.3333}}
  >{\raggedright\arraybackslash}p{(\linewidth - 4\tabcolsep) * \real{0.3333}}@{}}
\caption{Hipóteses pré-especificadas.}\tabularnewline
\toprule\noalign{}
\begin{minipage}[b]{\linewidth}\raggedright
ID
\end{minipage} & \begin{minipage}[b]{\linewidth}\raggedright
Hipótese
\end{minipage} & \begin{minipage}[b]{\linewidth}\raggedright
Comparação
\end{minipage} \\
\midrule\noalign{}
\endfirsthead
\toprule\noalign{}
\begin{minipage}[b]{\linewidth}\raggedright
ID
\end{minipage} & \begin{minipage}[b]{\linewidth}\raggedright
Hipótese
\end{minipage} & \begin{minipage}[b]{\linewidth}\raggedright
Comparação
\end{minipage} \\
\midrule\noalign{}
\endhead
\bottomrule\noalign{}
\endlastfoot
H1 & os atributos limpos preveem reúso entre donos & AUC do RF limpo
contra 0,5 \\
H2 & os atributos limpos reduzem o RMSLE de \texttt{owner\_count} & RF
limpo contra a média de \texttt{log1p} \\
H3a & a skill acrescenta a repositório e catálogo & limpo contra
repositório+catálogo (AUC) \\
H3b & repositório e catálogo acrescentam à skill & limpo contra skill
(AUC) \\
H4 & os atalhos acrescentam aos atributos limpos & limpo+atalhos contra
limpo (RMSLE) \\
\end{longtable}

A interpretação usa importância por permutação do RF limpo em dobras
agrupadas por dono, por coluna e por bloco (colunas do bloco permutadas
em conjunto). É associativa; atributos correlacionados compartilham
importância.

\subsection{Holdout novo}\label{holdout-novo}

Como o protocolo foi fixado após resultados exploratórios, a confirmação
exige dados nunca consultados. O holdout reúne conteúdos fora dos
200.000 amostrados, com \texttt{repository\_count\ \textgreater{}=\ 2},
cujos donos não aparecem entre as skills reutilizadas da amostra
(desenvolvimento e teste original), com limite de 50 por dono: 25.011
conteúdos de 10.484 donos, sem donos nem conteúdos em comum com a
amostra e com no máximo 37 conteúdos por dono. Viu-se apenas a
prevalência de \texttt{cross\_owner} (80,8\%) como verificação de
integridade. Os modelos escolhidos são treinados no desenvolvimento
inteiro e avaliados \textbf{uma única vez}, com IC por bootstrap de
donos (2.000 reamostras). O sucesso exige limite inferior do AUC acima
de 0,5 e limite superior da diferença de RMSLE para a baseline abaixo de
zero; uma trava impede reavaliações sem declaração explícita.
